\section{Dirichlet series with a witnessed margin}
\label{sec:dirichlet-interior}

Take $h\ge1$ and a complex argument $s$ with
$\mathop{\rm Re}s\ge1+1/h$. A statement about this variable in iPRA is a
schema for a specified primitive-recursive (PR) rational approximation to
$s$, with bounds on its size and approximation error. For rational complex
$s$ the resulting statement is uniform in its numerator and denominator
codes. A disk or half-plane condition refers to a supplied bound.
Write $u_n=\exp(-s\log n)$, $Z(s)=\sum_{n\ge1}u_n$. Then
\begin{align}
 D(s):=\sum_{n\ge1}\mu(n)u_n
    &= Z(s)^{-1}=\prod_{p\ {\rm prime}}(1-u_p),\label{eq:mob-ds}\\
 Q(s):=\sum_{n\ge1}|\mu(n)|u_n
    &=Z(s)/Z(2s).\label{eq:mob-sqdir}
\end{align}
The notation $Z$ agrees here with the usual Dirichlet-series definition of
$\zeta$. There are explicit finite error bounds for each equality.

\paragraph{Evaluation of the elementary weights.}
For $n\ge1$ put $v=(n-1)/(n+1)$. Use the series
$2\sum_{j\ge0}v^{2j+1}/(2j+1)$ for $\log n$. After the terms
$j<L$ its remainder is at most $2v^{2L+1}/(1-v^2)$.
If $0\le r\le1-1/b$ and $b\ge2$, then $r^b\le1/2$ by
$(1+1/b)^b\ge2$ and $1-1/b\le(1+1/b)^{-1}$.
Thus $b=\max(2,n+1)$ supplies a PR cutoff.
For $|z|\le B$, $B$ a positive integer, the exponential Taylor terms
from index $2B$ on have successive absolute ratio at most $1/2$,
and the term at $2B$ is at most $(1+B)^{2B}$.
For example, the tail from $2B+L$ on is bounded by
$2(1+B)^{2B}2^{-L}$. These estimates also justify replacing $s$ and
$\log n$ by sufficiently accurate rationals. On a ball of radius $B$
the exponential is Lipschitz with bound $3^B$ by its power series
and $\sum 1/j!\le3$.

For clarity, the algebra behind the weights is elementary too.
The binomial theorem in the product of the exponential series gives
$\exp(a+b)=\exp(a)\exp(b)$. As a formal identity in $v$,
\[
 \exp\!\left(2\sum_{j\ge0}\frac{v^{2j+1}}{2j+1}\right)
       =\frac{1+v}{1-v}.
\]
Indeed both sides have constant term $1$ and solve
$(1-v^2)y'=2y$, which determines the coefficients successively.
At $0\le v<1$ the nonnegative terms can be collected by degree:
truncate the inner sum using its geometric remainder, truncate the outer
exponential using the factorial remainder above, then take a degree cutoff
at the product of the two polynomial degrees. This proves
$\exp(\log n)=n$. On real arguments $\exp$ is strictly increasing with
a lower increment bound on each bounded interval, by
$\exp(t)\ge1+t$ for $t\ge0$ and the addition law. These facts give
$u_{ab}=u_a u_b$. Conjugation and the addition law give
$|\exp(z)|=\exp(\mathop{\rm Re}z)$.

\paragraph{One estimate for all the rearrangements.}
In $2^j\le n<2^{j+1}$ the sum of $|u_n|$ is at most
$2^{-\lfloor j/h\rfloor}$. Indeed there are $2^j$ terms and
$n^{1/h}\ge2^{\lfloor j/h\rfloor}$, as seen on taking the $h$th power.
Set $B_h=2h$. For $T=2^{hL}$ every finite sub-tail on $n\ge T$
has absolute mass at most
\[
 E_h(L)=2h\,2^{-L},\qquad \sum_{n\ge1}|u_n|\le B_h .
\]
For a finite upper cutoff each bound follows by grouping $h$ successive
blocks and summing a finite geometric progression. The same bounds hold
with coefficients $\mu$, $|\mu|$, or any coefficients of absolute value
at most $1$.

In the product $D(s)Z(s)$ restrict first to $ab\le T^2$.
The coefficient of $u_m$ is $\sum_{a\mid m}\mu(a)=[m=1]$.
Every discarded pair has $a>T$ or $b>T$; its total absolute contribution,
summed over both strips, is at most $2B_h E_h(L)$. This proves $D(s)Z(s)=1$.
In particular $|Z(s)|\ge1/B_h$, so division has a rational apartness
bound. Expand the finite product over $p\le T$ by choosing one summand
from each factor. Unique factorization identifies its terms with
$\mu(m)u_m$ over squarefree $m$ whose prime factors are all at most $T$.
Every missing term has $m>T$. Hence the product approximates $D(s)$ with
error at most $E_h(L)$. Grouping subsets by size also gives the
alternating sums over increasing prime tuples printed by MathWorld.

Finally, each $m$ has exactly one decomposition $m=a b^2$ with $a$
squarefree: split each prime exponent into its parity and twice its
integer half. In $Q(s)Z(2s)$ this gives coefficient $1$ at $u_m$.
Truncate at $ab^2\le T^3$; the discarded pairs lie in $a>T$ or $b>T$,
and $|u_{b^2}|\le |u_b|$. The same error bound proves
$Q(s)Z(2s)=Z(s)$. Applying the reciprocal identity at $2s$ permits
the division in \eqref{eq:mob-sqdir}.

\section{Lambert series and the exponential product}
\label{sec:mob-lambert}

These identities hold coefficient by coefficient in the truncated
polynomial rings $\mathbb Z[X]/(X^{H+1})$ or
$\mathbb Q[X]/(X^{H+1})$ as appropriate. They also hold as effective
complex identities on every disk $|q|\le \rho=1-1/b$, $b\ge2$:
\begin{align}
 \sum_{d\ge1}\frac{\mu(d)q^d}{1-q^d}&=q,\label{eq:mob-lam1}\\
 \sum_{d\ge1}\frac{\mu(pd)q^d}{q^d-1}
    &=\sum_{j\ge0}q^{p^j}\qquad(p\text{ prime}),\label{eq:mob-lam2}\\
 \prod_{d\ge1}(1-q^d)^{\mu(d)/d}&=\exp(-q).\label{eq:mob-logprod}
\end{align}
In the last identity the factor means
$\exp((\mu(d)/d)\log(1-q^d))$ with the logarithm given by its series
at $1$. This specifies its complex branch.

In \eqref{eq:mob-lam1} expand the geometric sum. The coefficient at
degree $m$ is $\sum_{d\mid m}\mu(d)=[m=1]$.
For \eqref{eq:mob-lam2} it is $-\sum_{d\mid m}\mu(pd)$.
Writing $m=p^a v$, $p\nmid v$, the summands with $p\mid d$ vanish,
while $-\mu(pd)=\mu(d)$ for $p\nmid d$. The answer is
$\sum_{d\mid v}\mu(d)=[v=1]$, including $m=1=p^0$.
For \eqref{eq:mob-logprod} add the logarithms:
\[
 -\sum_{d,j\ge1}\frac{\mu(d)}{dj}q^{dj}
 =-\sum_{m\ge1}\frac{q^m}{m}\sum_{d\mid m}\mu(d)=-q.
\]
Modulo $X^{H+1}$ all displayed sums are finite. For the product use
the truncated exponential $\sum_{i=0}^H y^i/i!$ at an argument
with zero constant term. Its addition law is the binomial theorem in
degrees at most $H$.

For the analytic readings the double sums are dominated by
$\sum_{d,j\ge1}\rho^{dj}$. The strip $d>K$ has mass at most
$\rho^{K+1}/(1-\rho)^2$ by summing first over $j$, and the strip
$j>K$ has the same bound. Thus cutting at total degree $dj\le K^2$
leaves error at most $2b^2\rho^{K+1}$. The outer-row tail alone is
at most $b^2\rho^{K+1}$. For target $2^{-t}$ one may take
$K=b(t+2b+2)$, since $\rho^b\le1/2$.
The lacunary sum uses the same degree cutoff and its remaining mass after
degree $K^2$ is at most $b\rho^{K^2+1}$.
In the logarithmic case each finite sum of logarithms has norm at most
$b^2$; exponentiating transfers the error with the factor $3^{b^2}$.
These estimates justify the exact rearrangements with only finite
geometric and factorial bounds.

\section{The three reciprocal-square masses}
\label{sec:mob-basel}

The values tabulated in MathWorld depend on two even zeta values. The
following finite polynomial proof supplies both and gives a PR rate.
We use angles in radians.

\paragraph{Effective trigonometry.}
Define sine and cosine by the imaginary and real parts of $\exp(it)$.
For $|v|<1$ put
$A(v)=\sum_{j\ge0}(-1)^jv^{2j+1}/(2j+1)$ and set
$\pi=4(A(1/2)+A(1/3))$. This is the usual arctangent normalization of
$\pi$. The alternating bounds give $2<\pi<4$.
The formal identity $\exp(2iA(v))=(1+iv)/(1-iv)$ follows by matching
the constant term and the coefficients of $(1+v^2)y'=2iy$.
At $|v|\le1/2$ absolute geometric and exponential majorants justify
evaluation as in Section~\ref{sec:dirichlet-interior}.
Multiplying the two rational values on the right for $v=1/2,1/3$
gives $i$. Thus $\exp(i\pi/2)=i$ and $\exp(i\pi)=-1$.

For $0\le x\le\pi/2<2$, the alternating series give
$x-x^3/6\le\sin x\le x$, in particular $\sin x\ge x/3$.
Also $\cos x=\sin(\pi/2-x)\ge0$ by the addition law, and
$|\sin x|,|\cos x|\le1$ by the complex norm identity.
We shall need $x\cos x\le\sin x$ on this interval. Indeed
$\sin x-x\cos x$ has successive signed terms
$(-1)^{j+1}2j x^{2j+1}/(2j+1)!$, $j\ge1$. Their absolute ratio is
$(j+1)x^2/[j(2j+3)(2j+2)]\le 2/5$ for $x\le2$.
Pairing successive terms proves the inequality. For $0<x<\pi/2$
we consequently have
\begin{equation}
 0\le x^{-2}-\cot^2 x\le1 .\label{eq:mob-cot-est}
\end{equation}
The upper bound is $x^{-2}\le1/\sin^2x=1+\cot^2x$, and the lower
bound follows from $x\cos x\le\sin x$.

\paragraph{A finite polynomial with known roots.}
Take $N=2m+1$ with $m\ge2$ and $x_k=k\pi/N$, $1\le k\le m$.
The polynomial $\mathop{\rm Im}(Y+i)^N$ (real indeterminate $Y$)
is an even polynomial in $Y$. Set $U=Y^2$ and divide by $N$ to obtain
the monic polynomial
\[
 P(U)=U^m-\frac{\binom N3}{N}U^{m-1}
           +\frac{\binom N5}{N}U^{m-2}+\cdots .
\]
For each $k$, $(\cot x_k+i)^N=(-1)^k/\sin^N x_k$, so
$r_k=\cot^2x_k$ is a root. There is explicit apartness between
different roots. For $x_j<x_k$,
$\cot x_j-\cot x_k=\sin(x_k-x_j)/(\sin x_j\sin x_k)
\ge \pi/(3N)$.
Further $\cot x_k\ge \cos x_k\ge(\pi/2-x_k)/3\ge\pi/(6N)$.
Thus differences of the $r_k$ have a positive rational lower bound.
Synthetic division successively extracts the $m$ linear factors:
evaluating a quotient at the next root gives zero because each previously
extracted factor there is apart from zero. This proves the coefficient
identities
\[
 \sum_k r_k=C,\qquad \sum_k r_k^2=C^2-2E,\qquad
 C=\frac{(N-1)(N-2)}6,\quad E=\frac{\binom N5}{N}.
\]

\paragraph{Passage to the sums.}
Set $t_k=(\pi/N)^2r_k$. By \eqref{eq:mob-cot-est},
$0\le k^{-2}-t_k\le (\pi/N)^2$. Their squares differ by at most
$2(\pi/N)^2$ since $t_k\le k^{-2}\le1$.
The tail of either $\sum k^{-2}$ or $\sum k^{-4}$ beyond $m$ is
at most $1/m$, using
$1/k^2\le1/(k-1)-1/k$.
Hence the respective sums differ from $(\pi/N)^2 C$ and
$(\pi/N)^4(C^2-2E)$ by at most $1/m+2m(\pi/N)^2$.
Now $C/N^2\to1/6$ and
$E/N^4=\prod_{j=1}^4(1-j/N)/120\to1/120$.
For an explicit check use
$0\le1-\prod_j(1-j/N)\le\sum_j j/N$ for $N\ge5$.
It gives $|C/N^2-1/6|\le1/(2N)$,
$|E/N^4-1/120|\le1/(12N)$ and
$|(C/N^2)^2-1/36|\le1/(6N)$.
All errors thus tend to zero with a PR modulus, e.g.\ taking
$m=128\cdot2^t$ makes each combined error less than $2^{-t}$.
We obtain $Z(2)=\pi^2/6$ and
$Z(4)=\pi^4(1/36-2/120)=\pi^4/90$.

At $s=2$ equations \eqref{eq:mob-ds}--\eqref{eq:mob-sqdir} now give
$D(2)=6/\pi^2$ and $Q(2)=15/\pi^2$. Since twice the indicator of
$\mu(n)=\pm1$ is $|\mu(n)|\pm\mu(n)$, the two separate masses are
\begin{equation}
 \sum_{\mu(n)=-1}\frac1{n^2}=\frac9{2\pi^2},\qquad
 \sum_{\mu(n)=1}\frac1{n^2}=\frac{21}{2\pi^2}.
\end{equation}
For all three positive series the rational partial sum through $m$
already has error at most $1/m$.

\paragraph{Lean implementation of the precision bounds.}
Represent a Gaussian rational by a pair of rational numerators with
positive denominators. For instance, the scalar cutoff
$\texttt{tailCut}(h,L)=2^{hL}$ and the disk cutoff
$\texttt{diskCut}(b,t)=b(t+2b+2)$ have type
$\mathbb N\to\mathbb N\to\mathbb N$. Define a finite series evaluation
from a list of coefficients, a cutoff and an input approximation.
The Dirichlet proof requires three finite lemmas: a bound for a dyadic
block, a bijection between divisor sums and hyperbola sums, and the
two-strip bound. An embedded derivation of the last lemma first scans the
rectangle by rows, includes all points with $a>T$, then includes only
$a\le T$ in the column strip. This uses a bounded sum of nonnegative
rational majorants. The supplied
\texttt{lean/FormalSeries.lean} verifies the ordinary Lambert coefficients
and the logarithmic coefficients of the product at every degree. It
represents $L$ by $l_n=n[X^n]L$ in positive degrees and represents $E$
by $e_n=n![X^n]E$. The formal equation $E'=L'E$ then reads
\[
 e_{n+1}=\sum_{k=0}^n n(n-1)\cdots(n-k+1)\,l_{k+1}e_{n-k},
\]
with an empty product at $k=0$. The file proves that $l_n=-[n=1]$
and that the normalized solution has $e_n=(-1)^n$.
For the remaining formal power series one may represent a
truncation by an array of $H+1$ rational codes. Cauchy multiplication
at index $k\le H$ is the bounded sum over $i\le k$. Formal differentiation
through degree $H$ reads an input through degree $H+1$.
Prove both coefficient ODE identities with that extra input coefficient.
The polynomial proof of the two zeta values uses coefficient descent in
synthetic division and the explicit lower bounds at the roots.
For real inequalities $A\le B$ with error radius $e$ on each input,
the target at a given precision is the rational comparison
$A_{\rm app}\le B_{\rm app}+2e$. Equality uses both inequalities.
The finite induction arguments can therefore be stated with a fixed precision
parameter and a bounded Boolean check for every earlier row, coefficient
or polynomial degree needed by the recursion.
Here is an error budget for inductions whose length depends on the
input. Generate the finite list of arithmetic operations and inference
steps for the chosen cutoffs; let $J$ be its length. Bound magnitudes
forward by replacing addition by the sum of bounds and multiplication
by their product. At each inverse or cancellation record the reciprocal
of the supplied lower bound. For example at the cotangent roots one may
use $|\sin x_k|\ge 2/(3N)$ and
$|r_j-r_k|\ge 4/(9N^2)$ for $j\ne k$.
Let $C\ge2$ be an integer exceeding all these bounds. With
$K=8(1+C)^4$, addition, multiplication, inversion on a ball of radius
less than $1/(2C)$ around an admissible denominator, and cancellation
against an apart factor all amplify incoming errors by at most $K$ per
step. Indeed a product error is at most $(2C+1)e$ and an inverse error
at most $2C^2e$ for input errors at most $e\le1/(2C)$.
Start every elementary numerical evaluation with error at most
$2^{-t}/(4CK^{J+1})$. Even with one new error of this size at each step,
the geometric sum bounds the final accumulated error by $2^{-t}$.
Power-series truncations are chosen by the preceding explicit tails.
All bounds, including the generation and length of the operation list,
use nested bounded loops or primitive recursion. This permits a single
quantifier-free induction on the list at the selected working precision,
including the synthetic-division argument with a variable number of roots.
